\section{Positive-time Gaussian limits and the uniform start}
\label{sec:gaussian-limits}

Two inputs give the positive-time edge evolution.  For bounded cylinder
sources, the finite Markov resolvent contraction bounds the optimizers
used in the stationary relaxed-energy theorem.  A variational argument
then gives the required strong semigroup convergence.  The warming
kernel bound controls the evolved densities, whose weak limits inherit
that evolution.  Warming also controls the total mass in a growing
block of edge modes.  A rotation argument distributes that mass among
individual modes and retains information specific to the uniform start.

Throughout this section we use Gaussian couplings and the completed
GOE matrix of Section~\ref{sec:background}, writing
\(\Lambda=\diag(\lambda_1,\ldots,\lambda_n)\).  Write
\(w_{n,s}=d(\nu_nP_{sT_n})/d\mu\) when the system size needs emphasis.
We use the represented-subsequence convention of that section and
denote the probability and expectation on a representation by
\(\widehat{\Prob}\) and \(\widehat{\E}\).

\begin{proposition}[Positive-time Gaussian subsequential limits]
\label{prop:gaussian-subsequence}
Every sequence of Gaussian system sizes tending to infinity has a further represented
subsequence on which the following assertions hold almost surely.  The
edge coordinates converge to \(g\), and there is a measurable random
family \(h_s\in L^2(\Pi_g)\), \(s>0\), with the following properties.
The family may at this stage depend on the represented data in addition
to \(g\).

There is a finite random \(C\ge1\), independent of \(s,t\), such that
\begin{equation}
 0\le h_s\le \mathcal K(s):=C\exp(C s^{-77/50}),\qquad
 \Pi_g[h_s]=1,\qquad
 h_{s+t}=\Sg_t h_s,\qquad s,t>0,
 \label{glim:density-family}
\end{equation}
where \(\Sg_t=\exp(-t c_*H_g)\).  The map \(s\mapsto h_s\) is continuous
in \(L^2(\Pi_g)\).  For each \(t>0\), \(\Sg_t\) has a nonnegative symmetric
kernel \(k_t^g\) relative to \(\Pi_g\) satisfying
\begin{equation}
 0\le k_t^g(x,y)\le \mathcal K(t)
 \quad\text{for \(\Pi_g\otimes\Pi_g\)-almost every \((x,y)\)}.
 \label{glim:limit-kernel}
\end{equation}
The constant \(77/50\) is the exponent \(1.54\) in
Theorem~\ref{thm:warming}.

On this subsequence,
\begin{equation}
 A_{n,J}\longrightarrow A_g
 \quad\text{in }C_{\mathrm{loc}}((0,\infty)),\qquad
 B_{n,J}\longrightarrow B_h
 \quad\text{in }C_{\mathrm{loc}}((0,\infty)^2),
 \label{glim:functional-subsequence}
\end{equation}
where
\begin{equation}
 B_h(s,t)=\sum_{a\ge1}\Pi_g[h_s x_a\,\Sg_t x_a].
 \label{glim:correlation-series}
\end{equation}
The series in \eqref{glim:correlation-series} converges locally uniformly
and is continuous.  Finally,
\begin{equation}
 \|h_s-1\|_{L^2(\Pi_g)}\longrightarrow0,\qquad
 B_h(s,\cdot)\longrightarrow A_g(\cdot)
 \quad\text{in }C_{\mathrm{loc}}((0,\infty))
 \quad(s\to\infty).
 \label{glim:relaxation}
\end{equation}
\end{proposition}

\begin{proof}
We first pass the semigroups on bounded sources, then construct the
densities, and finally justify the sum over all edge coordinates.

\paragraph{Varying \(L^2\) spaces.}
Choose a countable family \(\mathscr C\) of smooth compact cylinders
with this approximation property: every smooth compact cylinder
is approximated uniformly, together with its first derivatives, by
elements of \(\mathscr C\) supported in one fixed compact coordinate set.
Rational smooth approximations on rational boxes give such a family.
The stated approximation and the definition of the closed form make
\(\mathscr C\) dense in its form norm, and also dense in \(L^2(\Pi_g)\).
For linear tests put
\(\mathscr V=\operatorname{span}_{\mathbb Q}(\mathscr C\cup\{1\})\).
This is a countable family of bounded smooth cylinders, dense in
\(L^1(\Pi_g)\) and \(L^2(\Pi_g)\).  If \(\phi\in\mathscr V\) depends on the
first \(d\) coordinates, let
\(\phi_n=\phi(x_1,\ldots,x_d)\) on the finite cube.

The static cylinder and moment convergence of
\cite[Corollary~3.3]{FUA} gives
\begin{equation}
 \mu[\phi_n\psi_n]\longrightarrow\Pi_g[\phi\psi]
 \quad(\phi,\psi\in\mathscr V).
 \label{glim:static-pairings}
\end{equation}
On sequences bounded in \(L^2(\mu)\), the weak convergence of
Section~\ref{sec:background} can be tested on this countable family:
the pairings with every \(\phi\in\mathscr V\) converge to the
corresponding pairings with \(f\) if and only if the pairings with
every smooth compact cylinder do.  To pass from \(\mathscr V\) to an
arbitrary smooth compact cylinder, approximate it uniformly on a
common support by \(\mathscr C\), and use the uniform \(L^2\) bound
to control the pairing error in both spaces.  The converse follows
from the same approximation, with cutoffs tending to one for the
constant test.
A sequence bounded in \(L^2(\mu)\) has a weakly convergent subsequence:
diagonal extraction of its pairings defines a bounded linear functional
on the cylinder span by \eqref{glim:static-pairings}, and the Riesz
representation theorem supplies its \(L^2(\Pi_g)\) representative.
The same argument gives
\[
 \|f\|_{L^2(\Pi_g)}\le\liminf_n\|f_n\|_{L^2(\mu)}.
\]
As in Section~\ref{sec:background}, the convergence is strong when
the norms also converge.  For such a sequence,
\begin{equation}
 \lim_n\|f_n-\phi_n\|_{L^2(\mu)}^2
       =\|f-\phi\|_{L^2(\Pi_g)}^2
       \quad(\phi\in\mathscr C).
 \label{glim:common-approximation}
\end{equation}
Consequently strong convergence is equivalent to approximation, in both
spaces, by the same cylinders.  In particular, the pairing of an
\(L^2\)-bounded weakly convergent sequence with a strongly convergent
sequence converges.

\paragraph{Resolvents from bounded optimizers.}
Set
\[
 \mathfrak q_n=T_n\cE,\quad \mathfrak A_n=T_nH_n,\qquad
 \mathfrak q=c_*\Pi_g[|\nabla\cdot|^2],\quad \mathfrak A=c_*H_g.
\]
For \(\phi\in\mathscr C\) and rational \(z>0\), consider
\[
 u_n=(z+\mathfrak A_n)^{-1}\phi_n,\qquad
 u=(z+\mathfrak A)^{-1}\phi.
\]
The finite Markov property gives
\(\|u_n\|_\infty\le z^{-1}\|\phi\|_\infty\).
The resolvent identity also bounds \(\|u_n\|_2\) and
\(\mathfrak q_n(u_n)\) uniformly.  Thus these particular optimizers
satisfy the polynomial supremum bound as well as the norm and energy
bounds in the relaxed-energy theorem
\cite[Theorem~6.8]{FUA}.  That theorem applies to disorder-dependent
tests under the class-uniform quenched convention recorded in
Section~\ref{sec:background}.  This proof invokes it only for the
countably specified optimizer sequences, one for each
\(\phi\in\mathscr C\) and rational \(z>0\), on their common
represented event.  Its lower bound therefore gives
\begin{equation}
 \liminf_n\mathfrak q_n(u_n)\ge\mathfrak q(U)
 \quad\text{whenever \(u_n\) has weak limit \(U\)}.
 \label{glim:optimizer-liminf}
\end{equation}

Here is the variational passage, including the norm conclusion.  Let
\[
 v_n=\sup_f\{2\langle\phi_n,f\rangle_\mu-z\|f\|_2^2-\mathfrak q_n(f)\}
     =\langle\phi_n,u_n\rangle_\mu
     =z\|u_n\|_2^2+\mathfrak q_n(u_n),
\]
and define \(v\) by the analogous supremum in \(L^2(\Pi_g)\).  Take an
optimizer subsequence attaining the upper limit and then a weak limit
\(U\).  The norm lower bound and
\eqref{glim:optimizer-liminf} imply
\[
 \limsup_n v_n
 \le 2\langle\phi,U\rangle_{\Pi_g}
       -z\|U\|_2^2-\mathfrak q(U)
 \le v.
\]
For the reverse inequality, take \(F\in\mathscr C\).
The recovery part of \cite[Theorem~6.8]{FUA} supplies, for each fixed
recovery index \(j\), polynomially bounded finite functions strongly
converging to \(F\), whose energies converge to \(\mathfrak q(F)\) in
the order \(n\to\infty\), then \(j\to\infty\).  Testing the variational
supremum with these functions in exactly that order gives
\[
 \liminf_n v_n\ge
 2\langle\phi,F\rangle_{\Pi_g}-z\|F\|_2^2-\mathfrak q(F).
\]
The form density of \(\mathscr C\) allows the supremum over \(F\),
proving \(v_n\to v\).
Only the countably many sources in \(\mathscr C\) and rational \(z\)
are used on the common source event.
The limiting variational problem is strictly concave, so every weak
optimizer limit is its unique optimizer \(u\).  Moreover,
\[
 \limsup_n z\|u_n\|_2^2
 \le v-\mathfrak q(u)=z\|u\|_2^2.
\]
Together with norm lower semicontinuity this proves strong convergence
of \(u_n\) to \(u\).  This is the bounded-source version of the
variational argument in \cite[proof of Lemma~7.2]{FUA}.

The bound \(\|(z+\mathfrak A_n)^{-1}\|_{2\to2}\le z^{-1}\) and
\eqref{glim:common-approximation} now extend this resolvent convergence
to every strongly converging input.  No additional use of the lower
energy bound is needed for that extension.  To pass to semigroups, put
\(R_n=(1+\mathfrak A_n)^{-1}\), \(R=(1+\mathfrak A)^{-1}\), and, for
\(t>0\), define on \([0,1]\)
\[
 F_t(q)=
 \begin{cases}
 \exp[-t(q^{-1}-1)],&q>0,\\
 0,&q=0.
 \end{cases}
\]
This function is continuous.  Polynomial approximation of \(F_t\) and
convergence of every power of \(R_n\) on strongly converging inputs
prove
\begin{equation}
 P_{t_nT_n}f_n\longrightarrow \Sg_t f
 \quad\text{strongly whenever \(f_n\to f\) strongly and \(t_n\to t>0\)}.
 \label{glim:semigroup-strong}
\end{equation}
For the assertion with varying times, use
\[
 \sup_{\lambda\ge0}|e^{-t\lambda}-e^{-u\lambda}|
 \le\frac{|t-u|}{e\delta},\qquad t,u\ge\delta>0.
\]
This proof identifies the limit of \(P_{tT_n}\) as
\(\exp(-t c_*H_g)\) with the stated \(T_n=n^{2/3}\).

\paragraph{Densities and the limiting kernel.}
On the represented sequence, Theorem~\ref{thm:warming}, with fixed
admissible auxiliary parameters, gives a finite
random \(C\) and numbers \(\delta_n\downarrow0\) such that
\(\|P_{sT_n}\|_{1\to\infty}\le \mathcal K(s)\) for all \(s\ge\delta_n\).
Since \(w_{n,0}\) is a nonnegative density and the semigroup is
self-adjoint,
\begin{equation}
 0\le w_{n,s}\le \mathcal K(s),\qquad
 \mu[w_{n,s}]=1,\qquad \|w_{n,s}\|_2^2\le \mathcal K(s)
 \quad(s\ge\delta_n).
 \label{glim:finite-density}
\end{equation}
Include the pairings of these densities with
\(\mathscr V\) at each positive rational \(s\) in the
representation.  Every
limit functional \(L_s\) obeys
\[
 |L_s(\phi)|\le \mathcal K(s)\Pi_g[|\phi|],\qquad
 L_s(\phi)\ge-\mathcal K(s)\Pi_g[\phi^-],\qquad L_s(1)=1.
\]
Here \(\phi^-=\max\{-\phi,0\}\).  The bound follows from
\eqref{glim:finite-density} and static convergence also for the
continuous cylinders \(|\phi|\) and \(\phi^-\).
It extends \(L_s\) to \(L^1(\Pi_g)\).  The second inequality proves
positivity of the extension: if cylinders \(\phi_j\) converge in
\(L^1\) to \(f\ge0\), then
\(\Pi_g[\phi_j^-]\le\|\phi_j-f\|_1\to0\).
The extension is therefore represented by a density
\(0\le h_s\le \mathcal K(s)\) of mass one.  In particular
\(\|h_s\|_2^2\le \mathcal K(s)\).  This also justifies normalization if only
compact cylinders were initially retained: their cutoffs tending to
one have complements bounded by \(\mathcal K(s)\) times the corresponding
static probabilities.  Measurability can be seen concretely.  Take nested
finite cylinder partitions which generate the product Borel sets and
whose cell boundaries have zero \(\Pi_g\) measure; rational coordinate
boxes have this property by the cylinder densities of
\cite[Proposition~2.3]{FUA}.  The value \(L_s(\ind_A)\) for each cell is
a measurable limit of its values on fixed smooth approximations to
\(\ind_A\), which can in turn be approximated by \(\mathscr C\).
On each cell \(A\) put the ratio
\(L_s(\ind_A)/\Pi_g(A)\), with value zero if the denominator is zero.
These ratios are measurable, bounded by \(\mathcal K(s)\), and are the
conditional expectations of \(h_s\) on the nested partitions.  Their
bounded martingale limit gives a measurable representative of \(h_s\).

For positive rational \(s,t\), reversibility and
\eqref{glim:semigroup-strong} imply for every \(\phi\in\mathscr C\)
\[
 \Pi_g[h_{s+t}\phi]
 =\lim_n\mu[w_{n,s}P_{tT_n}\phi_n]
 =\Pi_g[h_s\Sg_t\phi].
\]
We used the weak--strong pairing rule; the \(w_{n,s}\) are weakly
convergent and bounded in \(L^2\).  Hence \(h_{s+t}=\Sg_t h_s\).
The closed cylinder-gradient form is a Dirichlet form: smooth normal
contractions obey its contraction inequality, and closure preserves it.
Constants belong to the form and have zero energy, by
\cite[Proposition~5.1]{FUA}.  Thus \(\Sg_t\) is Markov and preserves one.
For a real \(s>0\), choose a rational \(q<s\) and set
\(h_s=\Sg_{s-q}h_q\).  The rational relation makes this definition
independent of \(q\).  Strong continuity of the semigroup proves
\(L^2\) continuity in \(s>0\), as well as the relation for all \(s,t>0\).
Positivity and mass one persist.  Taking rationals \(q\uparrow s\) in
the density bound, using \(L^2\) convergence and an almost-everywhere
subsequence, proves \(h_s\le \mathcal K(s)\).  For deterministic dyadic rationals
\(q_j(s)\to s\), every cylinder integral at \(s\) is the limit of its
measurable values at \(q_j(s)\).  It is therefore jointly measurable
in the environment and \(s\).  The same partition construction
therefore gives jointly measurable representatives for the real-time
family.

The finite convergence against cylinders also holds along every
sequence \(s_n\to s>0\).  Indeed, for \(s\ge\delta>0\), the identity
\(w_{n,s}=P_{(s-\delta/2)T_n}w_{n,\delta/2}\) and spectral calculus give
\begin{equation}
 \|\partial_s w_{n,s}\|_2
 \le \frac{2}{e\delta}\sqrt{\mathcal K(\delta/2)}
 \quad(n\ \text{sufficiently large}).
 \label{glim:time-equicontinuity}
\end{equation}
Approximation of \(s\) by a rational time, followed by the weak
convergence there and the \(L^2\) continuity of \(h_s\), proves the
claim.

For completeness we obtain the kernel needed to sum coordinates.
For nonnegative smooth compact cylinders \(\phi,\psi\), the finite kernel
bound yields
\[
 0\le\mu[\phi_nP_{tT_n}\psi_n]
       \le \mathcal K(t)\mu[\phi_n]\mu[\psi_n].
\]
Passing to the limit using \eqref{glim:semigroup-strong}, and then
approximating indicators in \(L^2(\Pi_g)\), gives the same inequality
with \(\Pi_g[\phi\Sg_t\psi]\) in place of the left side.  Define on
rectangles the nonnegative functional
\[
 \ind_{A\times B}\longmapsto
       \Pi_g[\ind_A\Sg_t\ind_B].
\]
For a simple function on rectangles use common disjoint partitions
of the two factors.  The preceding rectangle inequality then bounds
the functional on every nonnegative such function by \(\mathcal K(t)\) times
its \(L^1(\Pi_g\otimes\Pi_g)\) norm.  It extends to a positive bounded
functional on that \(L^1\) space, so has a density \(k_t^g\) with
\(0\le k_t^g\le \mathcal K(t)\).  Self-adjointness makes the density symmetric.
Its marginals are \(\Pi_g\), since \(\Sg_t1=1\).  This proves
\eqref{glim:limit-kernel}; it does not require a joint pointwise version
of the kernels for all \(t\).

\paragraph{The tail of two equilibrium copies.}
Let \(x_a'=x_a(\sigma')\), and set
\[
 Q_{n,m}(\sigma,\sigma')=\sum_{a>m}x_a x_a',\qquad
 V_{n,m}=(\mu\otimes\mu)[Q_{n,m}^2].
\]
To see why this equilibrium quantity controls the quench trace, define
the relative transition kernel and the two-time density
\[
 p_{n,t}(\sigma,\sigma')
   =\frac{P_{tT_n}(\sigma,\sigma')}{\mu(\sigma')},\qquad
 d_{n,s,t}(\sigma,\sigma')
   =w_{n,s}(\sigma)p_{n,t}(\sigma,\sigma').
\]
The latter is the density of the two spins at times \(sT_n\) and
\((s+t)T_n\) relative to \(\mu\otimes\mu\).  On every compact subset
of \(s,t>0\), the warming bounds give a uniform bound on this density
for all sufficiently large \(n\) on the represented sequence.
Hence
\[
 \left|\sum_{a>m}\mu[w_{n,s}x_aP_{tT_n}x_a]\right|
 =\left|(\mu\otimes\mu)[d_{n,s,t}Q_{n,m}]\right|
 \le\|d_{n,s,t}\|_{L^2(\mu\otimes\mu)}V_{n,m}^{1/2}.
\]
The stationary trace has the same bound with \(d_{n,s,t}\) replaced
by \(p_{n,t}\).  Thus it suffices to control \(V_{n,m}\) uniformly
as \(m\) increases.
We claim that the representation may be further extracted so that
\begin{equation}
 \lim_{m\to\infty}\limsup_n V_{n,m}=0
 \quad\text{almost surely}.
 \label{glim:static-tail}
\end{equation}
Recall \(r_n=n^{-1/3+10^{-4}}\) from
Section~\ref{sec:edge-estimates}, and fix the spectrum in one of the
tight spectral sets of \cite[Proposition~2.3]{FUA}.  In this paragraph
take its canonical parameters
\[
 b_n^\circ=1+(-g_{1,n})_+,\qquad
 b^\circ=1+(-g_1)_+,\qquad
 \ell_{a,n}^\circ=(g_{a,n}+b_n^\circ)^{-1}.
\]
Let \(\mu^{\mathrm{sp}}\) be the spherical Gibbs law, and
\(L_\circ=Z_{\mathrm{cube}}/Z_{\mathrm{sp}}\), with both partition
functions relative to their normalized uniform measures.  Average
over the full Haar eigenframe at this fixed spectrum.  Multiplication
by \(L_\circ^2\) cancels the two cube Gibbs denominators.  In the
uncapped replica comparison of \cite[Lemma~3.2]{FUA}, the normalized
contribution of pairs with
\(|\rho|=|n^{-1}\sigma\cdot\sigma'|>C r_n\) is at most
\(C e^{-c n r_n^3}\).  Since \(|Q_{n,m}|^2\le n^{2/3}\), inserting
this square in the discarded part costs at most
\(C n^{2/3}e^{-c n r_n^3}\).

On the retained overlaps, the same comparison can be used with this
square inserted.  Here are the required uniform details.  Conditional
on \(\rho\), a Haar-rotated cube pair and a spherical pair have the
same law of an orthonormal sum-and-difference frame \((q,q')\).
Their spectral vectors can be written
\[
 Y(\rho)=\sqrt{\frac n2}
       \big(\sqrt{1+\rho}\,q+\sqrt{1-\rho}\,q'\big),\qquad
 Y'(\rho)=\sqrt{\frac n2}
       \big(\sqrt{1+\rho}\,q-\sqrt{1-\rho}\,q'\big).
\]
The ratio of the discrete overlap mass and the spherical mass on a
cell of width \(2/n\) is at most \(\exp(C+C n\rho^4)\), by the
replica comparison.  This is bounded on \(|\rho|\le C r_n\), because
\(n r_n^4=o(1)\).  The tilt exponent on this frame is
\[
 \frac n2\big((1+\rho)q^{\mathsf T}\Lambda q
                    +(1-\rho){q'}^{\mathsf T}\Lambda q'\big).
\]
Its oscillation in a cell is bounded when \(\|\Lambda\|\) is bounded.
If \(\rho,\widetilde\rho\) lie in that cell, the displayed vectors
change in Euclidean norm by at most \(C n^{-1/2}\).  After multiplying
by \(n^{-1/3}\), each vector changes by at most \(C n^{-5/6}\) and
has norm \(n^{1/6}\).  Projection onto the coordinates \(a>m\)
therefore gives, uniformly in \(m\),
\[
 |Q_{n,m}(\rho)-Q_{n,m}(\widetilde\rho)|\le C n^{-2/3}.
\]
The inequality \((u+v)^2\le2u^2+2v^2\) now bounds the inserted square
by the corresponding spherical square and \(C n^{-4/3}\), up to
the preceding bounded comparison factors.

Under two independent spherical Gibbs laws, sign symmetry kills the
cross terms in that square.  The uniform coordinate moment bound in
\cite[Proposition~2.3]{FUA} gives
\[
 (\mu^{\mathrm{sp}}\otimes\mu^{\mathrm{sp}})[Q_{n,m}^2]
 =\sum_{a>m}\big(\mu^{\mathrm{sp}}[x_a^2]\big)^2
 \le C\sum_{a>m}(\ell_{a,n}^\circ)^2.
\]
Since \(\ind_{\{L_\circ\ge1/2\}}\le4L_\circ^2\), the comparison just
proved yields
\begin{equation}
 \E_U[\ind_{\{L_\circ\ge1/2\}}V_{n,m}]
 \le C\sum_{a>m}(\ell_{a,n}^\circ)^2
       +C n^{-4/3}+C n^{2/3}e^{-c n r_n^3}.
 \label{glim:tail-haar}
\end{equation}
The constants are uniform on the retained spectral set.  The squared
resolvent tails tend to zero as \(m\to\infty\) there, and \(n r_n^3\)
is a positive power of \(n\).  Also
\(\Prob_U\{L_\circ<1/2\}=o(1)\), by
\cite[Lemma~2.5]{FUA}.  Markov's inequality, followed by exhaustion
of the tight spectral sets, proves for every \(\epsilon>0\)
\begin{equation}
 \lim_{m\to\infty}\limsup_n
       \Prob\{V_{n,m}>\epsilon\}=0.
 \label{glim:tail-probability}
\end{equation}
The same estimates make \(V_{n,m}\) tight for every fixed \(m\).
Include these countably many quantities in the representation and
extract coordinatewise limits \(V_m\).  The identity
\[
 V_{n,m}=\sum_{a,b>m}\mu[x_ax_b]^2
\]
shows that \(V_{n,m}\), and hence \(V_m\), decreases with \(m\).
Equation~\eqref{glim:tail-probability} forces \(V_m\downarrow0\)
almost surely.  Since \(V_{n,m}\to V_m\) for each fixed \(m\), this is
\eqref{glim:static-tail}.

The limiting version follows directly from the same sign symmetry
and moments.  With \(\ell_a^\circ=(g_a+b^\circ)^{-1}\), for \(M>m\),
\begin{equation}
 \left\|\sum_{m<a\le M}x_a x_a'\right\|_{L^2(\Pi_g\otimes\Pi_g)}^2
 =\sum_{m<a\le M}\Pi_g[x_a^2]^2
 \le C\sum_{a>m}(\ell_a^\circ)^2\longrightarrow0.
 \label{glim:limit-static-tail}
\end{equation}
Thus these interval sums have an \(L^2(\Pi_g\otimes\Pi_g)\) tail
whose norm tends to zero.

\paragraph{Correlation limits and relaxation.}
Orthogonality of the eigenframe gives the exact finite identities
\[
 A_{n,J}(t)=\sum_{a=1}^n\mu[x_aP_{tT_n}x_a],\qquad
 B_{n,J}(s,t)=\sum_{a=1}^n\mu[w_{n,s}x_aP_{tT_n}x_a].
\]
The preceding two-copy reduction and \eqref{glim:static-tail} give
uniform finite correlation tails on every compact positive time set.
For the limiting series, the corresponding two-time density is
\(h_s(x)k_t^g(x,x')\) relative to \(\Pi_g\otimes\Pi_g\);
for the stationary series it is \(k_t^g(x,x')\).
Equations~\eqref{glim:density-family} and~\eqref{glim:limit-kernel}
bound these densities uniformly on compact positive time sets.
The same Cauchy--Schwarz argument with
\eqref{glim:limit-static-tail} therefore gives the limiting uniform
tail control.

It remains to pass each fixed summand.  Static convergence of the mixed
pairings \(x_a\phi\) and the second moment \(x_a^2\) makes the finite
coordinate \(x_a\) strongly
convergent to the limiting \(x_a\).  If \(s_n\to s>0\), then
\(w_{n,s_n}x_a\) is weakly convergent to \(h_sx_a\).  To see this,
test first against a cylinder times a smooth cutoff of \(x_a\) at
height \(R\), and use the convergence after
\eqref{glim:time-equicontinuity}, extended from \(\mathscr C\) by the
common-support approximation above.  For a bounded cylinder \(\phi\),
the omitted part of the pairing is at most
\[
 \mathcal K(\delta)\|\phi\|_\infty
       \mu[|x_a|\ind_{\{|x_a|>R\}}]
 \le \frac{\mathcal K(\delta)\|\phi\|_\infty}{R}\mu[x_a^2]
 \quad(s_n\ge\delta).
\]
The fixed second moments are bounded and converge, so this error
vanishes uniformly as \(R\to\infty\).  Also
\(\|w_{n,s_n}x_a\|_2\le \mathcal K(\delta)\|x_a\|_2\).
The weak--strong pairing rule and
\eqref{glim:semigroup-strong} now give, whenever \(s_n\to s>0\) and
\(t_n\to t>0\),
\[
 \mu[w_{n,s_n}x_aP_{t_nT_n}x_a]
       \longrightarrow \Pi_g[h_sx_a\Sg_t x_a].
\]
The same argument without \(w_{n,s_n}\) gives the fixed stationary
summand.  Each limiting summand is continuous.  For \(t\ge\delta>0\), the kernel
bound gives
\[
 \|\Sg_t x_a\|_\infty\le \mathcal K(\delta)\Pi_g[|x_a|],\qquad
 \|x_a\Sg_t x_a\|_2\le \mathcal K(\delta)\|x_a\|_2^2.
\]
Thus \(L^2\) continuity of \(h_s\) gives continuity in \(s\),
uniformly for \(t\) in positive compact intervals.  Strong
continuity of \(\Sg_t x_a\), paired with \(h_sx_a\in L^2(\Pi_g)\),
gives continuity in \(t\).  The uniform tail bound proves the
asserted continuity and local uniform convergence of the full series.
It also proves the two locally uniform finite limits in
\eqref{glim:functional-subsequence}: failure of uniform convergence
on a compact set would provide a sequence of time points with a
convergent subsequence, contradicting the fixed-summand convergence,
the uniform tails, and continuity of the limiting sum.

Finally, the kernel of \(H_g\) consists exactly of the constants,
by \cite[Proposition~5.1]{FUA}.  The spectral theorem and
\(h_s=\Sg_{s-s_0}h_{s_0}\) give \(h_s\to\Pi_g[h_{s_0}]=1\) in
\(L^2(\Pi_g)\).  For each fixed \(a\), the last displayed bound
implies, uniformly for \(t\ge\delta\),
\[
 |\Pi_g[(h_s-1)x_a\Sg_t x_a]|
 \le \mathcal K(\delta)\|x_a\|_2^2\|h_s-1\|_2\longrightarrow0.
\]
The correlation tails are uniform also for \(s\ge s_0\), since
\(\|h_s\|_\infty\le \mathcal K(s_0)\).  Passing first through a finite sum and
then letting its length tend to infinity proves
\eqref{glim:relaxation}.
\end{proof}

The preceding proposition does not yet identify the family as a
function of \(g\).  The following estimate retains the information
from the finite uniform initialization that will make this
identification possible.

\begin{proposition}[An escape estimate from the uniform start]
\label{prop:uniform-seed-escape}
There is a measurable admissible choice \(b=b(g)\) such that, almost
surely,
\begin{equation}
 r_a:=g_a+b>0,\qquad D_g(b)>0,\qquad
 c_g(1+a)^{2/3}\le r_a\le C_g(1+a)^{2/3}\quad(a\ge1)
 \label{glim:rates}
\end{equation}
for finite positive \(c_g,C_g\).  Define
\begin{equation}
 M_g(x)=\frac12\sum_{a\ge1}\frac{x_a^2}{r_a}.
 \label{glim:weighted-size}
\end{equation}
Then \(0<M_g<\infty\) \(\Pi_g\)-almost surely.  Choose a Borel
full-measure set of edge realizations on which these assertions about
\(b\), \(r_a\), and \(M_g\) hold, with \(b\) Borel there and
\(r_a=g_a+b\).  On the product of this edge set with \(\R^\N\),
\(M_g\) is an extended nonnegative Borel function of \((g,x)\), since
its jointly Borel partial sums increase to it.  In every escape event,
\(x_i^2/M_g(x)\) denotes the quotient on the subset where
\(0<M_g(x)<\infty\), extended by zero to all remaining pairs in
\(\R^\N\times\R^\N\).  This version is jointly Borel and agrees with
the usual quotient \(\Pi_g\)-almost surely for every \(g\) in the
chosen edge set.

For every represented family in
Proposition~\ref{prop:gaussian-subsequence}, every fixed
\(i\ge1\), and every \(\epsilon>0\),
\begin{equation}
 \lim_{s\downarrow0}
 \widehat{\E}\!\left[
    \Pi_g\!\left[h_s\ind_{\{x_i^2/M_g>\epsilon\}}\right]\right]=0.
 \label{glim:seed-escape}
\end{equation}
Here the probability averages the represented environment as well as
the draw from \(h_s\Pi_g\).
\end{proposition}

\begin{proof}
We will show that \(x_i^2>s^{13/20}\) has vanishing averaged probability
under the limiting time-\(s\) law, while \(M_g<s^{14/25}\) also has
vanishing averaged probability as \(s\downarrow0\).  Away from these
two events, \(x_i^2/M_g\) is at most \(s^{9/100}\).
We obtain the numerator bound from
the finite uniform-start dynamics, then bound the denominator directly
under the conditioned edge measure.

\paragraph{The total mass in a block of finite modes.}
Fix \(\varepsilon>0\), and fix the auxiliary log-moment parameters
in Theorem~\ref{thm:warming} once and for all.  Its entropy bound and
Proposition~\ref{prop:edge-estimates} provide a quenched disorder event
\(\mathcal G_{n,\varepsilon}\) of limiting
probability at least \(1-\varepsilon\) on which their constants and
thresholds are fixed, the entropy bound holds simultaneously for
every positive time, and the uncapped cube projection tails hold at
every allowed scale.  Separately retain a spectrum-only event
\(\mathcal S_{n,\varepsilon}\) of
limiting probability at least \(1-\varepsilon\) for the band estimates
in Lemma~\ref{edge:spectral} and the tight leading-edge
bound \(|g_{1,n}|\le C_\varepsilon\).  The latter follows from the
GOE edge convergence in Section~\ref{sec:background}.  The constants in the
following calculations may depend on \(\varepsilon\).

For \(0<s<1\), put \(m=\lceil s^{-5/4}\rceil\).  For \(n\ge m\), set
\begin{equation}
 D_m(v)=\sum_{a\le m}x_a(v)^2,\qquad
 I=s^{47/100},\qquad R=s^{-53/100}.
 \label{glim:rotation-scales}
\end{equation}
Here \(x_a(v)\) is evaluated along the finite dynamics at physical
time \(vT_n\), and \(I=sR\).  Choose \(k_p=np^3=C' m\), with the
fixed \(C'\) sufficiently large.  Take \(s\) small enough that
\(C'm\) exceeds the fixed threshold, and then take \(n\) sufficiently
large that \(m\le n\) and \(p\le r_n\).  The spectral band estimates imply
that \(\dim\cH_p\ge m\) on \(\mathcal S_{n,\varepsilon}\), so \(P_p\)
contains the first \(m\) coordinates and
\(g_{m,n}\le T_np^2=(C'm)^{2/3}\).  Together with
\(|g_{1,n}|\le C_\varepsilon\), this gives on that event
\begin{equation}
 \max_{a\le m}|g_{a,n}|\le C m^{2/3}.
 \label{glim:finite-head-spectrum}
\end{equation}
At a fixed \(s\), the value
of \(k_p\) is fixed as \(n\to\infty\); this is why the extension to
the lowest scales in Proposition~\ref{prop:edge-estimates} is used here.

On \(\mathcal G_{n,\varepsilon}\cap\mathcal S_{n,\varepsilon}\) put
\(Y_p=\|P_p\sigma\|^2/M_p^2\).  Its
projection tail is
\(\mu\{Y_p>u\}\le C\exp(-a k_pu^3)\) for every \(u\ge u_0\), with
fixed \(u_0\).  Integration of this tail gives
\[
 \log\mu\exp\!\left(\tfrac a2 k_pY_p^3\right)\le C k_p.
\]
The entropy inequality for any probability density \(w\) relative
to \(\mu\), followed by Hölder's inequality, therefore gives
\[
 \mu[wY_p]\le
 C\left(1+\big(\Ent_\mu(w)/k_p\big)^{1/3}\right).
\]
Since \(D_m\le n^{-2/3}M_p^2Y_p=k_p^{1/3}Y_p\), the entropy part of
Theorem~\ref{thm:warming} yields under the direct dynamical law at
fixed disorder
\begin{equation}
 \E_{\nu_n}^J[D_m(v)]
 \le C\big(m^{1/3}+v^{-77/150}\big),\qquad 0<v\le s\le1.
 \label{glim:head-mass}
\end{equation}
This use of entropy is entirely relative to the ordinary Gibbs law
\(\mu\); the path in this expectation starts from \(\nu_n\).

Let \(Q(v)=\int_0^vD_m(w)\,dw\) and define the event
\[
 G=\{Q(s)<I,\ D_m(s)\le R\}.
\]
Integrating \eqref{glim:head-mass} and applying Markov's inequality
at the integral and the endpoint gives, on
\(\mathcal G_{n,\varepsilon}\cap\mathcal S_{n,\varepsilon}\),
\begin{equation}
 \Prob_{\nu_n}^J(G^c)
 \le C\big(s^{17/150}+s^{1/60}\big).
 \label{glim:cutoff-loss}
\end{equation}
Indeed \(m^{1/3}\le C s^{-5/12}\), and both comparisons use the
same two positive powers:
\[
 1-\frac5{12}-\frac{47}{100}
   =\frac{53}{100}-\frac5{12}=\frac{17}{150},\qquad
 1-\frac{77}{150}-\frac{47}{100}
   =\frac{53}{100}-\frac{77}{150}=\frac1{60}.
\]

\paragraph{A stopped rotation at fixed spectrum.}
For this step fix only a spectrum satisfying
\eqref{glim:finite-head-spectrum}.  Let \(\E_\Lambda\) average over
the full conditional Haar law of \(U\) and over the uniform-start
heat-bath path.  In particular, this expectation does not condition
on \(\mathcal G_{n,\varepsilon}\).

It is useful to specify the path likelihood.  In normalized time
\(v\), take independent site clocks of rate \(T_n\), a uniform
initial spin, and an independent fair mark
\(\eta\in\{-1,1\}\) at every ring, including rings at which the mark
equals the old spin.  This is a reference law; the marked path is
then determined without \(U\).  If the ring is at site \(k\), its
heat-bath probability is
\[
 p_{U,k}(\eta\mid\sigma)=
 \frac{e^{\eta b_k(\sigma)}}{2\cosh b_k(\sigma)},\qquad
 b_k(\sigma)=\sum_{\ell\ne k}W_{k\ell}\sigma_\ell.
\]
The likelihood of the dynamical path through \(s\) relative to the
reference law is
\begin{equation}
 \mathcal L_U=
 \prod_{\ell:v_\ell\le s}
 2p_{U,k_\ell}\big(\eta_\ell\mid\sigma(v_\ell-)\big).
 \label{glim:path-likelihood}
\end{equation}
There is no initial or clock factor depending on \(U\).
This likelihood fixes the marked spin outcomes while the matrix
varies; it permits differentiating the smooth heat-bath probabilities.
The absence of an initial score is the feature of the uniform start
used in this rotation.

For \(i,j\le m\), \(i\ne j\), right-rotate these two columns with
\(\dot u_i=u_j\) and \(\dot u_j=-u_i\).  Then
\[
 \dot x_i=x_j,\qquad \dot x_j=-x_i,\qquad
 \dot W=(\lambda_i-\lambda_j)
          (u_i u_j^{\mathsf T}+u_j u_i^{\mathsf T}).
\]
Writing \(\Delta_{ij}=\lambda_i-\lambda_j\), diagonal deletion in
the physical local field gives
\[
 \dot b_k=\Delta_{ij}\left[
  n^{1/3}(u_{j,k}x_i+u_{i,k}x_j)
                    -2u_{i,k}u_{j,k}\sigma_k\right].
\]
Orthogonality and
\(\sum_k u_{i,k}^2u_{j,k}^2\le1\), together with
\(|\Delta_{ij}|=T_n^{-1}|g_{i,n}-g_{j,n}|
\le C n^{-2/3}m^{2/3}\), imply
\begin{equation}
 \sum_k\dot b_k^2
 \le C n^{-4/3}m^{4/3}
       \big[n^{2/3}(x_i^2+x_j^2)+1\big].
 \label{glim:field-rotation}
\end{equation}
The final \(1\) is the contribution from deleting the diagonal.

In the filtration revealing \(U\) initially and then the marked
rings, the logarithmic derivative of \(\mathcal L_U\) through \(v\) is
\[
 S_{ij}(v)=\sum_{\ell:v_\ell\le v}
  \big(\eta_\ell-\tanh b_{k_\ell}(\sigma(v_\ell-))\big)
                 \dot b_{k_\ell}(\sigma(v_\ell-)).
\]
Under the dynamical law it is a square-integrable martingale.  In
fact each marked increment has conditional mean zero and conditional
variance \(\operatorname{sech}^2(b_k)\dot b_k^2\); its predictable
bracket is
\[
 \langle S_{ij}\rangle_v
 =T_n\int_0^v\sum_k\operatorname{sech}^2(b_k(\sigma(w-)))
                         \dot b_k(\sigma(w-))^2\,dw.
\]
Define
\[
 \tau=s\wedge\inf\{v:Q(v)\ge I\},\qquad
 I_a=\E_\Lambda\int_0^\tau x_a(v)^2\,dv\quad(a\le m).
\]
The function \(Q\) is continuous, so \(Q(\tau)\le I\).  Stopped
martingale isometry and \eqref{glim:field-rotation} give
\begin{equation}
 \E_\Lambda[S_{ij}(\tau)^2]
 \le C m^{4/3}\big(I_i+I_j+s n^{-2/3}\big).
 \label{glim:rotation-score}
\end{equation}

For a fixed marked reference history, \(D_m,Q,\tau\), and \(G\) are
exactly invariant under this rotation.  Apply Haar invariance first
to
\[
 F_{ij}=\int_0^\tau x_i(v)x_j(v)\,dv.
\]
Differentiating its integral with likelihood
\eqref{glim:path-likelihood}, and then projecting the terminal score
onto the stopped filtration, gives
\begin{equation}
 0=\E_\Lambda[\dot F_{ij}+F_{ij}S_{ij}(s)]
   =\E_\Lambda[\dot F_{ij}+F_{ij}S_{ij}(\tau)].
 \label{glim:haar-identity}
\end{equation}
Here is a justification at the stopping boundary.  The stopping
rule is constant on each rotation orbit of a fixed marked history,
so \(\dot F_{ij}\) differentiates only \(x_i x_j\).
At each fixed \(n\) and spectrum, the likelihood is bounded by
\(2^N\), where \(N\) is the number of rings, and its derivative is
bounded by this quantity times a constant (depending on \(n\) and
the spectrum) times \(N\).  Poisson exponential moments justify
differentiation of the full-time integral.  The random variable
\(F_{ij}\) is measurable at \(\tau\), and optional projection of
the square-integrable martingale replaces \(S_{ij}(s)\) by
\(S_{ij}(\tau)\).  This proves \eqref{glim:haar-identity} without
differentiating a law with a moving stopping time.

Since \(\dot x_i=x_j\) and \(\dot x_j=-x_i\), this identity reads
\[
 I_i-I_j=\E_\Lambda\!\left[
 S_{ij}(\tau)\int_0^\tau x_i(v)x_j(v)\,dv\right].
\]
Sum over \(j\le m\), \(j\ne i\).  The pathwise bound
\[
 \sum_{\substack{j\le m\\j\ne i}}\left(\int_0^\tau x_i x_j\,dv\right)^2
 \le \left(\int_0^\tau x_i^2\,dv\right)
       \left(\int_0^\tau D_m\,dv\right)
 \le I\int_0^\tau x_i^2\,dv
\]
and Cauchy--Schwarz on the product of probability and index spaces,
together with \(\sum_{j\le m} I_j\le I\) and
\eqref{glim:rotation-score}, yield
\begin{equation}
 m I_i\le I+
 \left\{C m^{4/3}
       (m I_i+I+m s n^{-2/3})\,I I_i\right\}^{1/2}.
 \label{glim:integrated-rotation}
\end{equation}
This estimate does not require independence between the scores.
To solve it, put
\[
 y=\frac{m I_i}{I},\qquad
 \alpha=C I m^{1/3},\qquad
 \kappa=\frac{m s n^{-2/3}}{I}.
\]
It gives \(y\le1+\sqrt{\alpha y(y+1+\kappa)}\).
Here \(I m^{1/3}=O(s^{4/75})\), while \(\kappa\to0\) as
\(n\to\infty\) at each fixed \(s\).  For sufficiently small \(s\)
and then sufficiently large \(n\), \(\kappa\le1\) and
\(\sqrt\alpha\le1/2\).  Using
\(\sqrt{y(y+2)}\le y+1\) proves
\begin{equation}
 I_i\le 3I/m\qquad(i\le m).
 \label{glim:individual-integral}
\end{equation}

Repeat \eqref{glim:haar-identity} with
\(F_{ij}=\ind_G x_i(s)x_j(s)\).  This random variable is measurable
at \(\tau\): it is zero if \(\tau<s\), and otherwise is measurable
at \(s=\tau\).  The indicator \(G\), including its boundary, is
constant on each rotation orbit, so its differentiation contributes
nothing.  Write \(Z_a=\E_\Lambda[\ind_G x_a(s)^2]\) for \(a\le m\).  Now
\(\sum_{a\le m} Z_a\le R\) and
\[
 \sum_{\substack{j\le m\\j\ne i}}\E_\Lambda[\ind_G x_i(s)^2x_j(s)^2]\le R Z_i.
\]
Use \eqref{glim:rotation-score} and
\eqref{glim:individual-integral} in the summed identity.  Taking \(n\)
large enough that \(m s n^{-2/3}\le I\), we obtain
\begin{equation}
 m Z_i\le R+\big(C m^{4/3} I R Z_i\big)^{1/2},
 \qquad Z_i\le C R/m.
 \label{glim:terminal-rotation}
\end{equation}
For the second inequality, \(z=mZ_i/R\) satisfies
\(z\le1+\sqrt{\alpha z}\), and hence \(z\le2+\alpha\).
As \(R/m\le C s^{18/25}\), Markov's inequality gives, at each of the
retained spectra,
\begin{equation}
 \Prob_\Lambda\{x_i(s)^2>s^{13/20},\,G\}
 \le s^{-13/20}Z_i\le C s^{7/100}.
 \label{glim:one-mode}
\end{equation}

We now combine the two probability laws.  Let \(\Prob_n\) include
the Gaussian disorder, the independent diagonal and eigenvector
signs, and the uniform-start dynamics.  Write \(\Prob_{\mathrm{dis}}\)
and \(\E_{\mathrm{dis}}\) for the completed disorder law and
\(\E_{\mathrm{eig}}\) for the eigenvalue law.  Then
\[
\begin{aligned}
 \Prob_n\{x_i(s)^2>s^{13/20}\}
 \le{}&
 \Prob_{\mathrm{dis}}\big((\mathcal G_{n,\varepsilon}
                         \cap\mathcal S_{n,\varepsilon})^c\big)\\
 &+\E_{\mathrm{dis}}\!\left[
   \ind_{\mathcal G_{n,\varepsilon}\cap\mathcal S_{n,\varepsilon}}
                   \Prob_{\nu_n}^J(G^c)\right]\\
 &+\E_{\mathrm{eig}}\!\left[
   \ind_{\mathcal S_{n,\varepsilon}}
      \Prob_\Lambda\{x_i(s)^2>s^{13/20},G\}\right].
\end{aligned}
\]
The second term uses \eqref{glim:cutoff-loss}; the last uses
\eqref{glim:one-mode} with its full conditional Haar average.
The complements of
\(\mathcal G_{n,\varepsilon}\) and \(\mathcal S_{n,\varepsilon}\)
cost at most \(2\varepsilon\) in the limit, so for
each fixed \(i\) and all sufficiently small \(s\),
\begin{equation}
 \limsup_n\Prob_n\{x_i(s)^2>s^{13/20}\}
 \le2\varepsilon+
 C_\varepsilon\big(s^{17/150}+s^{1/60}+s^{7/100}\big).
 \label{glim:finite-numerator}
\end{equation}
This separation is essential: no frame-dependent quenched event
was imposed inside the Haar identity.  Sending \(s\downarrow0\)
and then \(\varepsilon\downarrow0\) proves the finite numerator
estimate.

For a fixed positive \(s\), the cylinder convergence in
Proposition~\ref{prop:gaussian-subsequence} transfers this estimate
to \(\rho_s(dx)=h_s(x)\Pi_g(dx)\).  Indeed each coordinate marginal
of \(\Pi_g\) has a density by
\cite[Proposition~2.3]{FUA}; the same is true of \(\rho_s\).
The boundary \(x_i^2=s^{13/20}\) thus has zero limiting probability.
Weak convergence of this coordinate marginal and bounded convergence
over the represented environment yield
\begin{equation}
 \lim_{s\downarrow0}
       \widehat{\E}\big[\rho_s\{x_i^2>s^{13/20}\}\big]=0.
 \label{glim:limit-numerator}
\end{equation}
Only finite-coordinate convergence was used in this step.

\paragraph{The weighted denominator.}
Corollary~\ref{edge:limit-growth} gives the two-sided bound in
\eqref{glim:rates} for every admissible \(b\), with measurable constants
when \(b\) is measurable.  It remains to choose \(b\) with \(D_g(b)>0\).
Start with a measurable admissible \(b_0\), for example
\(b_0=1+(-g_1)_+\).  The resolvent-difference identity in
\cite[Proposition~2.3]{FUA} gives, for \(b\ge b_0\),
\[
 D_g(b)-D_g(b_0)=
 \sum_a\left(\frac1{g_a+b_0}-\frac1{g_a+b}\right).
\]
The right side increases to \(+\infty\) as \(b\to\infty\), because
\(\sum_a(g_a+b_0)^{-1}=\infty\).  Choose the first integer
increment \(b=b_0+j\) for which \(D_g(b)>0\).  This is measurable
and admissible, so Corollary~\ref{edge:limit-growth} gives
\eqref{glim:rates}.

Put \(\ell_a=r_a^{-1}\).  The uniform coordinate moment bound in
\cite[Proposition~2.3]{FUA} gives
\[
 \Pi_g[M_g]\le C_g\sum_a\ell_a^2<\infty.
\]
The first coordinate has no atom at zero, so \(M_g>0\) almost surely.
This proves that \eqref{glim:weighted-size} is well defined under
\(\Pi_g\) and under every \(\rho_s\).

Let \(N=\lceil s^{-8/5}\rceil\) and consider the coordinates
\(N\le a\le2N\).  Under the independent Gaussian law of the
conditioning construction, write \(Z_a=\sqrt{\ell_a}\xi_a\) with
independent standard normals \(\xi_a\).  Equation~\eqref{glim:rates}
implies
\[
 \frac12\sum_{a=N}^{2N}\ell_a Z_a^2
 \ge c_gN^{-4/3}\sum_{a=N}^{2N}\xi_a^2.
\]
The chi-square lower-tail bound shows that this is at least
\(c_gN^{-1/3}\), except on an event of probability at most
\(C_g e^{-c_gN}\).

The same estimate survives the norm conditioning with a fixed cost.
Leave a sufficiently large fixed head of \(m_0\) coordinates outside
this block.  Let \(q\) be the density of
\(\sum_a(Z_a^2-\ell_a)\), and \(q_{m_0}^{\rm head}\) the density of its
first \(m_0\) terms.  The tail disintegration in
\cite[Proposition~2.3]{FUA} states that
\begin{equation}
 \frac{d(\Pi_g)_{>m_0}}{d\bigotimes_{a>m_0}N(0,\ell_a)}
 =
 \frac{q_{m_0}^{\rm head}\!\left(
 D_g(b)-\sum_{a>m_0}(x_a^2-\ell_a)\right)}{q(D_g(b))}.
 \label{glim:conditioned-tail-density}
\end{equation}
The head density is bounded by Fourier inversion once \(m_0\) is
large enough, and \(q(D_g(b))>0\).  Hence the ratio in
\eqref{glim:conditioned-tail-density} is bounded by a finite
\(g\)-dependent constant.  For all sufficiently small \(s\), the
block lies outside the head and
\[
 c_g N^{-1/3}\ge s^{14/25},
 \qquad\text{since}\qquad
 \frac{14}{25}-\frac8{15}=\frac2{75}>0.
\]
It follows that
\begin{equation}
 \Pi_g\{M_g<s^{14/25}\}\le C_g e^{-c_g s^{-8/5}}.
 \label{glim:denominator-static}
\end{equation}
Using \eqref{glim:density-family}, for almost every represented
environment we obtain
\[
 \rho_s\{M_g<s^{14/25}\}
 \le C_g C\exp\!\left(C s^{-77/50}-c_gs^{-8/5}\right)
 \longrightarrow0,
\]
because \(8/5>77/50\).  These probabilities are bounded by one, so
the convergence also holds after \(\widehat{\E}\).

On the complement of the two exceptional events in
\eqref{glim:limit-numerator} and \eqref{glim:denominator-static},
\[
 \frac{x_i^2}{M_g(x)}
 \le s^{13/20-14/25}=s^{9/100}.
\]
This tends to zero and proves \eqref{glim:seed-escape}.
\end{proof}

For each represented limiting law, the averaged limit in
\eqref{glim:seed-escape} yields a deterministic rational sequence
\(s_k\downarrow0\) along which, for almost every represented
environment, the displayed conditional probabilities tend to zero
for every fixed \(i\) and \(\epsilon>0\).
Section~\ref{sec:entrance-selection} proves this extraction and uses
the resulting criterion.
